\chapter{Working from the command line}\label{ch:console}

Two scripts in the repository's root solve projects without opening a window. They do not
need PySide6.

\section{run\_headless.py}

\file{run\_headless.py} runs the study cases of any project, prints every result and can
export what it ran:

\begin{terminal}
python run_headless.py examples/lf_220kv.psn                   # every case of the project
python run_headless.py examples/lf_220kv.psn --case LFd        # one case (repeatable)
python run_headless.py examples/s2_grid.psn --study thevenin   # a study with its defaults
python run_headless.py examples/lf_220kv.psn --set A=0         # change a variable first
python run_headless.py examples/s2_grid.psn --export s2.json   # also write a file
python run_headless.py examples/s2_grid.psn --list             # cases and studies
\end{terminal}

\begin{reftable}{@{}>{\raggedright\arraybackslash}p{3.6cm}X@{}}
\toprule
\sffamily\bfseries Option & \sffamily\bfseries What it does \\ \midrule
\cli{project.psn} & The project to solve. \\
\cli{--case ID} & Run this case of the project. Repeatable. Without \cli{--case} or
  \cli{--study}, every case runs, or every study with its defaults if the project has no
  cases. \\
\cli{--study KEY} & Run this study: the project's case of the same id if there is one,
  otherwise the study with its defaults. Repeatable. The keys are \code{model},
  \code{s1a}, \code{s1b}, \code{s1c}, \code{ybus}, \code{linear}, \code{thevenin},
  \code{short\_circuit}, \code{loadflow}, \code{seqmodel} and \code{unbalanced\_load}. \\
\cli{--set NAME=VALUE} & Set a project variable before solving. Repeatable. \cli{--set A=0}
  gives the answers printed in the assignments. \\
\cli{--export PATH} & Also write the results to a file. The extension (\file{.txt},
  \file{.csv}, \file{.tex}, \file{.json}) picks the format (\cref{sec:export-results}). \\
\cli{--quiet} & Do not print the results; useful with \cli{--export}. \\
\cli{--list} & List the project's cases (id, study, title) and the available studies,
  and stop. \\
\cli{--run-script} & Run the variable script the project links to. Without it, the values
  saved in the project are used. \\
\bottomrule
\end{reftable}

The output starts with the project's title, then the variables set and a note if a linked
script was not run, then each result as in the \gui{Log} tab.

\paragraph{Exit codes.} \code{0} when every case ran. \code{1} when at least one case
failed: its reason is printed and the other cases still run. \code{2} when the project
cannot be opened, a \cli{--case} names a case that does not exist, a \cli{--set} is
malformed or its expression fails, or the export file has an unknown extension. This
makes the runner usable in scripts and batch checks:

\begin{terminal}
python run_headless.py examples/lf_220kv.psn --case LFa --set A=0 --quiet --export lf_a0.json
\end{terminal}

\section{run\_s1\_headless.py}

\file{run\_s1\_headless.py} prints the S1 report: bases, per-unit values, S1a, S1b and
S1c. It runs on the built-in S1 network or on a project.

\begin{terminal}
python run_s1_headless.py                             # built-in S1 network, A = 80
python run_s1_headless.py --a-number 0                # the answer published in the statement
python run_s1_headless.py examples/s1_a80.psn         # a saved project
python run_s1_headless.py my_network.psn --study s1b  # a single study (repeatable)
\end{terminal}

\begin{reftable}{@{}>{\raggedright\arraybackslash}p{3.6cm}X@{}}
\toprule
\sffamily\bfseries Option & \sffamily\bfseries What it does \\ \midrule
\cli{project.psn} & The project; without it, the built-in S1 network. \\
\cli{--a-number N} & The A-number of the built-in network (80 by default). Ignored for a
  project. \\
\cli{--study KEY} & Only this study, any key. Repeatable. By default \code{model},
  \code{s1a}, \code{s1b} and \code{s1c}. \\
\cli{--run-script} & Run the project's linked variable script. \\
\bottomrule
\end{reftable}

Each study uses the options of the project's case of the same id, so the ports and buses
chosen in the window apply. It prints the network's validation messages first and
returns exit code 2 when the project cannot be opened or the studies cannot run.

\section{From Python}

The engine is a plain Python package (\code{psa.core}) with no Qt in it, so a script can
use it directly:

\begin{terminal}
from psa.core.serialize import load
from psa.core.study import run_case
from psa.core.export import export_results

proj = load("examples/lf_220kv.psn")
proj.params.set_var("A", "0")
proj.apply_parameters()
results = [run_case(proj.network, case, proj.params) for case in proj.cases]
print(results[0].scalars["P_loss_MW"].value)       # 26.424...
export_results(results, "lf.json")
\end{terminal}

A \code{Result} holds the formatted tables and the numbers behind them: \code{scalars},
per-bus \code{buses}, per-element \code{elements} and labelled \code{matrices}. These are
the same names as in the JSON export.
